Soit $f$ la fonction dérivable sur $\R$ définie par \quad
$f(x) = (3x- 1)\e^{-x} +2$.

On note $\mathcal{C}_f$ sa courbe représentative dans un repère orthonormé \Oij.


\begin{enumerate}
	\item Déterminer la limite de la fonction $f$ en $- \infty$.
	\item Déterminer la limite de la fonction $f$ en $+ \infty$.

	Interpréter graphiquement ce résultat.
	\item On note $f'$ la fonction dérivée de $f$.
	\begin{enumerate}
		\item Vérifier que, pour tout nombre réel $x,\quad f'(x) = (4 - 3x)\e^{-x}$.
		\item Étudier le signe de la dérivée et en déduire le tableau de variations de la fonction $f$ sur $\R$.
	\end{enumerate}
	\item \begin{enumerate}
		\item Démontrer que l'équation $f(x) = 0$ admet une unique solution sur $\R$.

		On notera $\alpha$ cette solution.

		\item Donner une valeur approchée de $\alpha$ à $10^{-2}$  près.
		\item En déduire le tableau de signes de la fonction $f$ sur $\R$.
		\item Soit la fonction $F$ définie sur $\R$ par : $F(x) =(- 3x - 2)\e^{-x} + 2x$.

		On admet que la fonction $F$ est une primitive de $f$.

		Calculer, en unités d'aire, l'aire du domaine limité par la courbe $\mathcal{C}_f$, l'axe de
		abscisses et les droites d'équation $x = 0$ et $x = 2$.
	\end{enumerate}
	\item  On admet que la droite d'équation $y = 4x + 1$ est la tangente à la courbe $\mathcal{C}_f$ au point d'abscisse 0.

	Vérifier que, pour tout réel $x,\: (3x - 1)\e^{-x} + 2 \leqslant 4x + 1$.
\end{enumerate}
